\section{Local Protection Competition and the Main Result}\label{sec:main}

The preceding section isolates two responses to public protection. First, a safer jurisdiction becomes more attractive as a primary production location. Second, lower primary risk reduces firms' willingness to maintain geographically separate backup capacity. This section shows that the interaction of those responses can reverse the social policy ranking.

\subsection{The two marginal incentives}

At a symmetric policy profile, define the planner's marginal payoff from a common increase in public protection by
\begin{equation}
M_S(\gamma)\equiv \left.\frac{dW(a,a)}{da}\right|_{a=0},
\label{eq:MS}
\end{equation}
and define jurisdiction $A$'s marginal payoff from a unilateral increase by
\begin{equation}
M_L(\gamma)\equiv \left.\frac{\partial G_A(a_A,0)}{\partial a_A}\right|_{a_A=0}.
\label{eq:ML}
\end{equation}
The planner internalizes the change in product-market surplus and private redundancy. A local government internalizes only its assigned share of national market surplus but additionally values the increase in its expected primary-plant count.

The wedge can be written without using the canonical witness. Let
\begin{equation}
\lambda_0\equiv
\left.\frac{\partial p(a_A,a_B)}{\partial a_A}\right|_{(0,0)}
\label{eq:lambda0}
\end{equation}
be the plant-attraction response at the symmetric origin.

\begin{proposition}[Plant-attraction threshold]\label{prop:attraction-threshold}
At the symmetric origin, under the quadratic protection cost in \eqref{eq:national-welfare},
\begin{equation}
M_L=\frac14 M_S+2b\lambda_0.
\label{eq:marginal-decomposition}
\end{equation}
Hence if $\lambda_0>0$ and $M_S<0$, the local marginal incentive is positive if and only if
\begin{equation}
b>\frac{-M_S}{8\lambda_0}.
\label{eq:attraction-threshold}
\end{equation}
\end{proposition}

\begin{proof}
At $(0,0)$, the marginal resource cost of protection is zero and symmetry implies $\mathcal S_{a_A}=\mathcal S_{a_B}$. Hence
\[
M_S=2\mathcal S_{a_A}.
\]
From equation~\eqref{eq:government-payoff},
\[
M_L=\frac12\mathcal S_{a_A}+2b\lambda_0.
\]
Substituting $\mathcal S_{a_A}=M_S/2$ gives equation~\eqref{eq:marginal-decomposition}. If $\lambda_0>0$ and $M_S<0$, then $M_L>0$ is equivalent to
\[
2b\lambda_0>-\frac14 M_S,
\]
or, equivalently, equation~\eqref{eq:attraction-threshold}.
\end{proof}

\begin{remark}[General incidence and hosting benefits]\label{rem:general-incidence}
The equal-share and linear-hosting specification in \eqref{eq:government-payoff} is used for transparency rather than for the marginal mechanism. For the local comparison only, suppose
\[
\widetilde G_A=\omega\mathcal S+B(2p)-C(a_A),
\]
where $\omega>0$ is a fixed local weight on national market surplus, $B$ is differentiable at the symmetric expected plant count $2p=1$, and $C'(0)=0$. Let the coordinated benchmark value the same hosting benefits and policy costs,
\[
\widetilde W
=
\mathcal S+B(2p)+B\!\left(2(1-p)\right)-C(a_A)-C(a_B).
\]
Along the common-policy path, symmetry gives $p(a,a)=1/2$. Hence the hosting terms are constant on that path and, at the origin,
\[
\widetilde M_S=2\mathcal S_{a_A},
\qquad
\widetilde M_L
=
\omega\mathcal S_{a_A}+2B'(1)\lambda_0
=
\frac{\omega}{2}\widetilde M_S+2B'(1)\lambda_0.
\]
Thus, when $\lambda_0>0$ and $\widetilde M_S<0$, the marginal Protection--Attraction Conflict obtains whenever
\[
B'(1)>\frac{-\omega\widetilde M_S}{4\lambda_0}.
\]
The baseline follows from $\omega=1/2$, $B(n)=bn$, and $C(a)=ca^2/2$. This remark concerns the local sign decomposition only; the global planner optimum and global best-response statements in Theorem~\ref{thm:main} continue to rely on the maintained baseline specification and its exact certificates.
\end{remark}

The decomposition separates the national resilience margin from the jurisdictional attraction motive. Their signs depend on the downstream resilience and location game. If the social marginal can be written as $M_S=E_0+R_0$, where $E_0>0$ is the direct engineering effect of reducing primary risk while holding private redundancy fixed and $R_0<0$ is the induced private-redundancy response, then a marginal Protection--Attraction Conflict obtains whenever
\begin{equation}
-R_0>E_0
\qquad\text{and}\qquad
8b\lambda_0>-E_0-R_0.
\label{eq:transparent-conflict}
\end{equation}
The first inequality says that redundancy crowd-out dominates the direct engineering benefit at the national margin; the second says that the local plant-attraction benefit overturns one quarter of that national marginal loss. These inequalities characterize the marginal sign conflict $M_S<0<M_L$ under the stated decomposition. The global equilibrium result below additionally requires planner optimality over the policy square and government best responses over each full policy interval.

For the rational parameter vector
\begin{equation}
\theta_0=
\left(
\gamma,q_0,k,H,b,c,\bar a
\right)
=
\left(
\frac{12}{25},\frac3{10},\frac{169}{3000},\frac1{100},\frac{11}{200},3,\frac2{25}
\right),
\label{eq:witness}
\end{equation}
exact rational evaluation gives
\begin{equation}
M_S(12/25)<0<M_L(12/25).
\label{eq:marginal-conflict}
\end{equation}
Thus, at the same policy profile, coordinated welfare falls when both jurisdictions add protection, while a jurisdiction gains from adding protection unilaterally.

\begin{theorem}[Protection--Attraction Conflict]\label{thm:main}
Within the maintained symmetric primitive family described in Section~\ref{sec:model}, there exists a nonempty open set of primitives for which coordinated welfare is maximized by zero additional local climate protection while the decentralized jurisdictional game has a strictly positive symmetric protection equilibrium. At the canonical rational witness in \eqref{eq:witness}, the coordinated optimum is uniquely $(a_A,a_B)=(0,0)$ and the positive symmetric decentralized equilibrium satisfies
\begin{equation}
a^{LG}\in(0.02301968485,\;0.02301968495).
\label{eq:alg-interval}
\end{equation}
Every downstream continuation used to evaluate first-stage deviations is unique and interior on the admissible policy domain $[0,\bar a]^2$.
\end{theorem}

The proof is computer-assisted but exact at the global-policy step. Exact rational root isolation certifies the unique symmetric local-government first-order-condition root. Tensor-product Bernstein certificates over the full policy rectangle prove that coordinated welfare is strictly decreasing in each protection coordinate, backup continuations remain regular and interior, the private-information location game is globally unclipped and has a unique interior probability fixed point, and the local government's objective is strictly concave in its own policy when the rival chooses the isolated equilibrium root. Hence the isolated stationary point is a global best response and the positive symmetric profile is a Nash equilibrium. Numerical global-deviation routines are retained only as independent stress tests. Appendix~\ref{app:verification} states the supporting lemmas, the exact certificate principle, and the logical implication from the certified inequalities to Theorem~\ref{thm:main}.

\subsection{A product-substitutability interval}

Holding $(q_0,k,H,b,c,\bar a)$ at the witness values and varying $\gamma$, exact root isolation identifies a local-government threshold
\begin{equation}
\gamma_L\in(0.4745893334177,\;0.4745893334186)
\end{equation}
and a planner threshold
\begin{equation}
\gamma_S\in(0.4982912217039,\;0.4982912217048).
\end{equation}
Hence $\gamma_L<\gamma_S$. For $\gamma$ between these thresholds, the local marginal incentive at zero is positive while the coordinated marginal incentive is negative. Figure~\ref{fig:policy-regime} plots these \emph{origin marginals}, evaluated by exact rational differentiation before conversion for plotting. These thresholds do not identify global policy equilibria, planner optima, or their boundaries over the displayed interval. The global policy ranking is proved at the canonical witness and on its open neighborhood, as stated in Theorem~\ref{thm:main}.

\begin{figure}[t]
\centering
\includegraphics[width=.78\textwidth]{../figures/policy_regime.pdf}
\refstepcounter{figure}
\label{fig:policy-regime}
\par\smallskip\textbf{Fig.~\thefigure}\quad Origin marginal incentives
\end{figure}

\subsection{Economic interpretation}

The conflict concerns the industrial value of place-based protection. Safer primary production reduces the private value of maintaining costly geographic backup capacity. The induced readiness loss can make additional common protection undesirable even after engineering benefits and preparedness-cost savings are counted. A local government can still gain because the same protection attracts primary production from its rival. In Proposition~\ref{prop:attraction-threshold}, the attraction term $2b\lambda_0$ offsets one quarter of the national marginal loss. The theorem establishes a region in which these incentives support positive decentralized protection while coordinated industrial welfare is maximized at zero.
